\documentclass[11pt,a4paper]{article}
\usepackage[margin=1in]{geometry}
\usepackage{amsmath,amssymb,amsthm}
\usepackage{algorithm}
\usepackage{algpseudocode}
\usepackage{booktabs}
\usepackage{hyperref}

\theoremstyle{definition}
\newtheorem{definition}{\textbf{Definition}}[section]
\newtheorem{theorem}{\textbf{Theorem}}[section]
\newtheorem{lemma}[theorem]{\textbf{Lemma}}
\newtheorem{remark}{\textbf{Remark}}[section]

\title{\textbf{Hierarchical Memory Offloading: CPU RAM and NVMe Streaming for Trillion-Parameter Scale}}
\author{\textbf{Team Scaibu} \\ \small Email: \texttt{jaydeep.9664920749@gmail.com} \\ \small \textbf{Architecture and Core Engine Team}}
\date{\textbf{October 2026}}

\begin{document}
\maketitle

\section{Definitions and Cognitive Mental Models}

\subsection{Formal Mathematical Definition}
\textbf{Definition (Hierarchical Memory Offloading and Parameter Swapping):}
When model scale $P$ exceeds the aggregate VRAM capacity of all GPUs in a cluster, \textbf{Memory Offloading} (Rajbhandari et al., 2021 - ZeRO-Infinity; Ren et al., 2021 - ZeRO-Offload) extends the memory hierarchy by evicting inactive states from fast GPU HBM ($1\text{--}3\,\text{TB/s}$) to host CPU DRAM ($30\text{--}60\,\text{GB/s}$ via PCIe Gen4/Gen5) and non-volatile NVMe storage ($3\text{--}7\,\text{GB/s}$).

Memory states are categorized into three distinct offloading regimes:
\begin{enumerate}
    \item \textbf{Optimizer State Offloading}:
    Adam optimizer states ($12P$ bytes in FP32) are permanently stored in host CPU memory. Gradients computed on the GPU are streamed over PCIe to the host CPU, which updates the master weights and streams the updated FP16/BF16 weights back to the GPU.
    \item \textbf{Activation Offloading / Checkpointing}:
    Forward activations $h_l$ are copied to CPU DRAM during the forward pass and freed from GPU VRAM. During the backward pass, $h_l$ is streamed back to the GPU immediately prior to layer $l$'s backward execution.
    \item \textbf{Parameter Offloading (ZeRO-Infinity)}:
    All $L$ transformer layer weights reside on NVMe SSD or CPU RAM. Layer $l$'s weights $W_l$ are prefetched asynchronously into GPU VRAM immediately before execution, and evicted immediately upon layer completion:
    {\boldmath
    \begin{equation}
    \mathcal{M}_{\text{GPU}}^{\text{resident}} = \mathcal{M}_{\text{layer}} + \mathcal{M}_{\text{act\_local}} \ll \mathcal{M}_{\text{total}}
    \end{equation}
    }
\end{enumerate}

\subsection{Expanded Non-Technical Definition (Intuition for Anyone)}
\begin{enumerate}
    \item \textbf{The Desk, the Filing Cabinet, and the Basement}:
    \begin{itemize}
        \item \textbf{GPU VRAM}: Your small office desk. Ultra-fast, but only fits 1 open binder at a time.
        \item \textbf{CPU RAM}: The filing cabinet next to the desk. Takes 3 seconds to reach into, but holds 50 binders.
        \item \textbf{NVMe SSD}: The basement archive. Takes 30 seconds to fetch from, but holds 1,000 binders.
    \end{itemize}
    \item \textbf{Streaming Just-In-Time}: You only need chapter 5 on your desk while you are reading chapter 5. While your eyes read chapter 5, an assistant is already walking to the filing cabinet to grab chapter 6!
    \item \textbf{Trillion-Parameter Training on a Single Server}: Without offloading, training a 500-billion parameter model requires dozens of supercomputers. With CPU/NVMe offloading, you can fit the entire model on a single server, trading runtime speed for immense financial savings.
\end{enumerate}

\subsection{Child-Friendly Mental Model (Physical Metaphor)}
Imagine an artist painting a 100-color mural with a tiny palette:
\begin{itemize}
    \item \textbf{Monolithic VRAM}: Trying to fit all 100 paint buckets directly on your tiny wooden hand palette. You drop the palette and spill everything!
    \item \textbf{Offloading}: Keeping the 100 paint buckets on the floor. You dip your brush into Color 1, paint, wash the brush, and dip into Color 2. Your hand palette never runs out of space!
\end{itemize}

\section{Core Mathematical Equations and Definitions}

\subsection{PCIe Layer Transfer Latency}
{\boldmath
\begin{equation}
t_{\text{swap}} = \frac{M_{\text{layer}}}{\mathcal{B}_{\text{PCIe}}}
\end{equation}
}
\begin{itemize}
    \item \textbf{Formal Definition}: Unidirectional transfer latency for layer weights $M_{\text{layer}}$ over bandwidth $\mathcal{B}_{\text{PCIe}}$.
    \item \textbf{What It Means}: Time required to stream layer parameters from host memory into GPU VRAM.
    \item \textbf{Why It Is Important}: Must be less than or equal to forward execution time to achieve zero-overhead prefetching.
\end{itemize}

\subsection{Effective Throughput Penalty}
{\boldmath
\begin{equation}
\mathcal{P}_{\text{throughput}} = \frac{t_{\text{comp}} + \max(0, \, t_{\text{swap}} - t_{\text{comp}})}{t_{\text{comp}}}
\end{equation}
}
\begin{itemize}
    \item \textbf{Formal Definition}: Ratio of offloaded step duration to ideal non-offloaded step duration.
    \item \textbf{What It Means}: Measures the exact slowdown introduced when PCIe bandwidth cannot keep up with compute.
    \item \textbf{Why It Is Important}: Governs whether CPU or NVMe offloading is economically viable.
\end{itemize}

\section{Rigorous Mathematical Analysis Checklist}

\begin{itemize}
    \item \textbf{Notation}: $M_{\text{layer}}$ (layer size in MB), $\mathcal{B}_{\text{PCIe}}$ (bandwidth in GB/s), $L$ (total layers).
    \item \textbf{Epistemic Status}: Rajbhandari et al. (ZeRO-Infinity, ISCA 2021). Foundation of extreme parameter scale.
    \item \textbf{Dependency Graph}: Host Memory $\to \text{Async Stream Layer } l+1 \to \text{Compute Layer } l \to \text{Evict Layer } l \to \text{Next}$.
    \item \textbf{Dimensions}: Model layers $L \in [32, 128]$, PCIe bandwidth $\in [32, 128]$ GB/s.
    \item \textbf{Regimes}: Hardware-constrained enterprise environments, single-node fine-tuning of 70B+ models.
    \item \textbf{Invariants}: GPU memory utilization never exceeds physical capacity $\mathcal{M}_{\text{GPU}}$.
    \item \textbf{Isomorphisms}: Isomorphic to virtual memory paging and demand paging in modern operating systems.
\end{itemize}

\section{Theoretical Theorems and Formal Proofs}

\begin{theorem}[Zero-Overhead Memory Offloading Condition]
Let $t_{\text{comp}}(l)$ be the computation execution time of layer $l$, and let $M_{\text{layer}}(l+1)$ be the parameter size of layer $l+1$. If the interconnect bandwidth satisfies:
{\boldmath
\begin{equation}
\mathcal{B}_{\text{interconnect}} \ge \frac{M_{\text{layer}}(l+1)}{t_{\text{comp}}(l)}
\end{equation}
}
then asynchronous prefetching completely hides memory transfer latency behind computation, incurring zero throughput penalty ($\mathcal{P}_{\text{throughput}} \equiv 1.0$).
\end{theorem}

\begin{proof}
Let asynchronous streams be double-buffered such that layer $l+1$'s parameters begin streaming from CPU DRAM over PCIe at timestep $\tau_l^{\text{start}}$, which coincides with the initiation of layer $l$'s GEMM kernel.
The parameter transfer completes at:
{\boldmath
\begin{equation}
\tau_{\text{transfer\_done}} = \tau_l^{\text{start}} + \frac{M_{\text{layer}}(l+1)}{\mathcal{B}_{\text{interconnect}}}
\end{equation}
}
The computation of layer $l$ completes at:
{\boldmath
\begin{equation}
\tau_{\text{comp\_done}} = \tau_l^{\text{start}} + t_{\text{comp}}(l)
\end{equation}
}
Layer $l+1$ can begin computation without idle stalling if and only if $\tau_{\text{transfer\_done}} \le \tau_{\text{comp\_done}}$:
{\boldmath
\begin{equation}
\tau_l^{\text{start}} + \frac{M_{\text{layer}}(l+1)}{\mathcal{B}_{\text{interconnect}}} \le \tau_l^{\text{start}} + t_{\text{comp}}(l) \implies \mathcal{B}_{\text{interconnect}} \ge \frac{M_{\text{layer}}(l+1)}{t_{\text{comp}}(l)}
\end{equation}
}
When this condition holds, the GPU stream switches to layer $l+1$ with zero idle barrier cycles, proving that memory capacity can be expanded arbitrarily into CPU RAM with zero performance loss.
\end{proof}

\begin{theorem}[Upper Bound on Active Resident GPU Memory]
Under layer-wise parameter offloading with activation recomputation, the peak memory resident on the GPU is strictly bounded by the maximum single-layer footprint:
{\boldmath
\begin{equation}
\mathcal{M}_{\text{peak}} \le 2 \max_{l} M_{\text{layer}}(l) + \max_l \mathcal{M}_{\text{act}}(l)
\end{equation}
}
independent of the total network depth $L$.
\end{theorem}

\begin{proof}
Under double-buffered prefetching:
1. The GPU holds the parameters of active layer $l$ ($M_{\text{layer}}(l)$) while concurrently buffering incoming prefetch parameters for layer $l+1$ ($M_{\text{layer}}(l+1)$).
2. The parameter memory is at most $M_{\text{layer}}(l) + M_{\text{layer}}(l+1) \le 2 \max_l M_{\text{layer}}(l)$.
3. Under activation recomputation (checkpointing), forward activations from previous layers are discarded immediately upon layer completion, so only the local layer's working activation buffer $\mathcal{M}_{\text{act}}(l)$ resides in memory.
Summing both components:
{\boldmath
\begin{equation}
\mathcal{M}_{\text{peak}} \le 2 \max_l M_{\text{layer}}(l) + \max_l \mathcal{M}_{\text{act}}(l)
\end{equation}
}
Since this upper bound contains no dependence on total layer count $L$, a model with 1,000 layers consumes no more GPU memory than a model with 10 layers.
\end{proof}

\section{Foundational Architectural Questions and Epistemic Meta-Questions}

\subsection{Architectural Question 1: CPU Adam vs. GPU Adam}
\textbf{Question:} Does running Adam updates on the host CPU cause training bottlenecks?\\
\textbf{Answer:} Modern high-core count server CPUs (AMD EPYC, Intel Xeon) with AVX-512 can execute FP32 Adam updates in tens of milliseconds. In ZeRO-Offload, CPU Adam updates run concurrently with GPU operations across microbatches, maintaining over $80\%$ of pure GPU training speed while eliminating $75\%$ of GPU memory requirements.

\subsection{Architectural Question 2: NVMe SSD Wear and Durability}
\textbf{Question:} Does streaming weights to and from NVMe SSDs destroy enterprise SSD drive endurance?\\
\textbf{Answer:} Reading weights from NVMe generates zero SSD write wear. Only optimizer state writes and periodic checkpointing write to the drive. With modern enterprise U.2/U.3 PCIe Gen4 SSDs rated for 1–3 Drive Writes Per Day (DWPD), continuous offloaded training can run for years without drive degradation.

\section{Algorithmic Implementation Specification}

\begin{algorithm}[H]
\caption{Hierarchical Memory Offloading and Swap Feasibility Evaluation}
\begin{algorithmic}[1]
\Require GPU VRAM capacity $C_{\text{vram}}$, total layers $L$, layer params $M_{\text{param}}$, layer acts $M_{\text{act}}$, bandwidth $\mathcal{B}$, target
\Ensure Resident VRAM $M_{\text{gpu}}$, offloaded RAM $M_{\text{host}}$, swap latency $t_{\text{swap}}$, is\_offloaded, total size $M_{\text{tot}}$
\State $M_{\text{tot}} \gets L \cdot (M_{\text{param}} + M_{\text{act}})$
\State $\text{is\_offloaded} \gets (M_{\text{tot}} > C_{\text{vram}})$
\If{$\text{is\_offloaded}$}
    \State $M_{\text{gpu}} \gets M_{\text{param}} + M_{\text{act}}$
    \State $M_{\text{host}} \gets M_{\text{tot}} - M_{\text{gpu}}$
    \State $B_{\text{eff}} \gets (\text{target} == \text{"nvme"}) \; ? \; 0.25 \cdot \mathcal{B} \cdot 1024.0 : \mathcal{B} \cdot 1024.0$
    \State $t_{\text{swap}} \gets (M_{\text{param}} / B_{\text{eff}}) \cdot 1000.0$
\Else
    \State $M_{\text{gpu}} \gets M_{\text{tot}}, \quad M_{\text{host}} \gets 0.0, \quad t_{\text{swap}} \gets 0.0$
\EndIf
\State \Return $(M_{\text{gpu}}, M_{\text{host}}, t_{\text{swap}}, \text{is\_offloaded}, M_{\text{tot}})$
\end{algorithmic}
\end{algorithm}

\section{Big-$\Theta$ Complexity Analysis}

\begin{itemize}
    \item \textbf{Offload Memory Evaluation Time}: $\Theta(1)$ constant operations.
    \item \textbf{PCIe Streaming Transfer Volume}: $\Theta(L \cdot M_{\text{param}})$ bytes per training step.
    \item \textbf{Peak GPU Memory Utilization}: $\Theta(M_{\text{layer}})$ invariant to model depth.
    \item \textbf{Host Memory Footprint}: $\Theta(L \cdot M_{\text{layer}})$ scalable into terabytes of system RAM.
\end{itemize}

\section{Single Canonical Repository Reference}

The complete deterministic scalar implementation, verification suite, and unit test benchmarks are published at:
\begin{center}
\url{https://github.com/Chief-Strategist-J/llm-obs-infra}
\end{center}

\end{document}
